% Fictional examples only. Create local/cards.tex for private or submission-specific cards.

\createCard{A1.7}
  {Für das fiktive Lernteam von Beispielorganisation sind gemeinsame Arbeitsregeln für Aufgabenübergaben, Besprechungen und Reviews erarbeitet, in zwei Workshops abgestimmt und in einer einseitigen Teamvereinbarung festgehalten worden.}
  {Die Vereinbarung wird von allen vier Beispielrollen verwendet. Neue Aufgaben enthalten dadurch immer eine zuständige Person, einen Review-Schritt und ein vereinbartes Ergebnis.}
  {Fiktives Lernteam in Beispielorganisation}
  {0}
  {0}
  {1}
  {Projektmanagement, Leadership}

\createCard{A2.4}
  {Das Ergebnis eines fiktiven Projekts für eine digitale Pflanzenbibliothek ist mit einer kurzen Demo und drei auf die Zielgruppe abgestimmten Beispielszenarien präsentiert worden.}
  {Die Präsentation dauerte zwölf Minuten. Alle drei Szenarien konnten ohne zusätzliche Erklärung nachvollzogen und die zwei vereinbarten Rückfragen direkt beantwortet werden.}
  {Fiktives Lernprojekt, Beispielorganisation}
  {1}
  {0}
  {0}
  {Präsentationstechnik, Schriftliche Kommunikation}

\createCard{A3.2}
  {Grundlagen zu barrierearmen Webtexten sind mit einem Lernplan, kurzen Übungsaufgaben und einer anschliessenden Überarbeitung von fünf Beispielseiten erschlossen worden.}
  {Die fünf Beispielseiten erfüllen die zuvor festgelegte Checkliste. Die verwendeten Lernnotizen und die Checkliste sind als nachvollziehbare Arbeitsgrundlage abgelegt.}
  {Fiktive Weiterbildung bei Beispielorganisation}
  {0}
  {1}
  {0}
  {Selbst- und Kompetenzmanagement, Web Engineering}

\createCard{B7.3}
  {Die technischen Anforderungen für einen fiktiven Ausleihkalender sind aus drei erfundenen Nutzungsszenarien abgeleitet, priorisiert und in einer klaren Spezifikation mit Akzeptanzkriterien festgehalten worden.}
  {Jede der zwölf Anforderungen ist einer Rolle, einem Szenario und mindestens einem prüfbaren Akzeptanzkriterium zugeordnet. Offene Punkte sind separat markiert und nicht als erledigt ausgegeben worden.}
  {Fiktives Konzeptprojekt in Beispielorganisation}
  {1}
  {0}
  {0}
  {Software Engineering, Projektmanagement}
